\section*{第 \ref{ch:b1:mammal-organization} 章——哺乳动物的功能组织}

\begin{solution}{exo:b1:mammal-organization:1}
脊椎动物：两侧对称，\omterm{def:b1:mammal-organization:bodyplan}{头部集中化}，背侧有包围脊髓的脊椎轴，
腹侧有容纳内脏的\omterm{def:b1:mammal-organization:bodyplan}{体腔}，四肢。
哺乳动物：\omterm{def:b1:mammal-organization:bodyplan}{膈肌}将\omterm{def:b1:mammal-organization:bodyplan}{体腔}分隔为胸腔和腹腔，
有毛发，有乳汁，体温受调节。
\end{solution}

\begin{solution}{exo:b1:mammal-organization:2}
营养：消化（胃）、呼吸（肺）、循环（心脏）、排泄（肾脏）。
联系：神经（脑）、感觉\omterm{def:b1:mammal-organization:organ}{器官}（眼睛）、内分泌（甲状腺）、
骨骼和肌肉（股骨、二头肌）、皮肤。
生殖：生殖\omterm{def:b1:mammal-organization:organ}{器官}（卵巢、睾丸）、乳腺。
\end{solution}

\begin{solution}{exo:b1:mammal-organization:3}
水：$0.6\times 60 = \qty{36}{L}$；\omterm{def:b1:mammal-organization:compartments}{细胞内液} \qty{24}{L}；
\omterm{def:b1:mammal-organization:compartments}{细胞外液} \qty{12}{L}，其中\omterm{def:b1:mammal-organization:compartments}{组织液} \qty{9}{L}，\omterm{def:b1:mammal-organization:compartments}{血浆}
\qty{3}{L}。
\end{solution}

\begin{solution}{exo:b1:mammal-organization:4}
$1.1/5 = 22\%$。\qty{1.1}{L/min} 即每天 \qty{1580}{L}，
相当于血液总量的 290 倍。
\end{solution}

\begin{solution}{exo:b1:mammal-organization:5}
剩余菊粉 \qty{1.7}{g}；$V = 1.7/0.12 = \qty{14}{L}$
为\omterm{def:b1:mammal-organization:compartments}{细胞外液}；\omterm{def:b1:mammal-organization:compartments}{细胞内液} $42 - 14 = \qty{28}{L}$。
\end{solution}

\begin{solution}{exo:b1:mammal-organization:6}
毛细血管壁可让水和小离子自由通过，故二者达到平衡，
但阻挡蛋白质。\omterm{def:b1:mammal-organization:compartments}{血浆}蛋白以渗透作用将水吸入血管；
若没有它们，水将渗入
组织间隙，使其肿胀（水肿）；
\omterm{def:b1:mammal-organization:compartments}{血浆}体积也将下降。
\end{solution}

\begin{solution}{exo:b1:mammal-organization:7}
并联时，每个\omterm{def:b1:mammal-organization:organ}{器官}都接受充足动脉血氧与全额血压，
其血流可独立调节。串联时，脑所接受的血液已被肌肉耗用，
压强较低，
故其供血恰在运动时下降。
\end{solution}

\begin{solution}{exo:b1:mammal-organization:8}
$hS\Delta T = 80 - 20 = \qty{60}{W}$; $h = 60/(1.8\times 13) =
\qty{2.6}{W.m^{-2}.K^{-1}}$。当空气温度为 \qty{33}{\celsius} 时，
非蒸发散热消失；皮肤温度必须高于气温，
血管扩张可使其升至接近 \qty{36}{\celsius}，
其余热量则须由出汗带走。
\end{solution}

\begin{solution}{exo:b1:mammal-organization:9}
需散去的热量为 $800 - 200 - 150 = \qty{450}{W}$，即 \qty{1620}{kJ/h}，
每小时出汗 \qty{675}{g}。若不出汗，$450/(70\times 3500) =
\qty{1.8e-3}{K/s}$，约 \qty{0.11}{K/min}：半小时内升温 \qty{3}{K}。
\end{solution}

\begin{solution}{exo:b1:mammal-organization:10}
小型动物单位质量散热最快（$S/V$ 大），
需要持续、高效、可靠的产热：棕色脂肪无需运动即可
以稳定速率化学产热；而战栗需要发达的肌肉（新生儿缺乏），
会干扰运动与摄食，且是间歇性的。
棕色脂肪还靠近大血管，
可直接加热血液。
\end{solution}

\begin{solution}{exo:b1:mammal-organization:11}
升温的下丘脑感知到“过热”：狗尽管处于寒冷中仍喘息，
皮肤血管扩张，核心温度在一小时内下降一两度。
若没有下丘脑，便不存在\omterm{def:b1:organism-environment:homeostasis}{调定点}：狗在寒冷中既不战栗，
在炎热中也不喘息，
其体温在两种情况下都会趋向室温。
\end{solution}

\begin{solution}{exo:b1:mammal-organization:12}
\omterm{def:b1:mammal-organization:compartments}{细胞外液}（\omterm{def:b1:mammal-organization:compartments}{组织液}和\omterm{def:b1:mammal-organization:compartments}{血浆}）的成分（富含钠）、温度和体积几乎恒定，
如同海水浸润最早的动物那样，浸润着每一个细胞。
其恒定性由\omterm{def:b1:organism-environment:homeostasis}{负反馈}回路维持：
肾脏调节体积和离子，肺调节气体，
肝脏和胰腺调节葡萄糖，下丘脑调节温度。
其“海岸”是交换表面——肠道、肺、肾脏、皮肤——
即它与外界相遇之处；
循环则是搅动它的水流。
\end{solution}

\begin{solution}{pb:b1:mammal-organization:1}
\textbf{1.} $B = 3.4\times 2^{0.75} = \qty{5.7}{W}$.
\textbf{2.} 洞穴：$5.7\times 12\times 3600 = \qty{246}{kJ}$；
觅食：$17.1\times 6\times 3600 = \qty{369}{kJ}$；洞外：
$11.4\times 6\times 3600 = \qty{246}{kJ}$.
\textbf{3.} \qty{861}{kJ}，即 $861/493 = 1.75$ 倍基础代谢
（$B$ 在 \qty{24}{h} 内 $= \qty{493}{kJ}$）。
\textbf{4.} $0.25\times 18\times 0.65 = \qty{2.9}{kJ/g}$.
\textbf{5.} $861/2.9 = \qty{295}{g}$ 鲜草，
含 \qty{74}{g} 干物质。
\textbf{6.} $295/2000 = 15\%$，占其体重。
\textbf{7.} $861/20 = \qty{43}{L}$ 的 $\mathrm{O_2}$；\qty{30}{mL/min}。
\textbf{8.} $0.9\times 43 = \qty{39}{L}$ 的 $\mathrm{CO_2}$，即 $39/22.4\times 44 = \qty{76}{g}$。
\textbf{9.} $0.75\times 295 = \qty{221}{g}$.
\textbf{10.} 消化的干物质 $0.65\times 74 = \qty{48}{g}$；水
$0.6\times 48 = \qty{29}{g}$.
\textbf{11.} 未消化物 \qty{26}{g}；粪便 $26/0.6 = \qty{43}{g}$，
含水 \qty{17}{g}。
\textbf{12.} $(44 - 5)\times 0.86 = \qty{34}{g}$.
\textbf{13.} 输入：$221 + 29 = \qty{250}{g}$。输出：$17 + 130 + 34 + 20
= \qty{201}{g}$。盈余 \qty{49}{g}。
\textbf{14.} 不会：草提供的物质多于其损失，盈余
以额外尿液排出（共约 \qty{180}{mL}）。
\textbf{15.} 相同的干物质 \qty{74}{g} 存在于 $\qty{148}{g}$ 的草中；
食物中的水为 \qty{74}{g}；输入 \qty{103}{g} 对输出 \qty{201}{g}，
每天亏损 \qty{100}{g}，
须饮水或从体内补足。
\textbf{16.} 水 $0.6\times 2.0 = \qty{1.2}{L}$；\omterm{def:b1:mammal-organization:compartments}{细胞外液}
\qty{0.4}{L}；\omterm{def:b1:mammal-organization:compartments}{血浆} \qty{0.1}{L}。
\textbf{17.} $1.5/15 = \qty{0.10}{L}$：一致。
\textbf{18.} $0.10/0.60 = \qty{0.167}{L}$，占体重的 \qty{8}{\%}。
\textbf{19.} $5\times(2/70)^{0.75} = \qty{0.35}{L/min}$；循环时间
$0.167/0.35 = \qty{0.48}{min}$，约 \qty{29}{s}。
\textbf{20.} $350/200 = \qty{1.75}{mL}$，每次搏动。
\textbf{21.} $0.1\times 2^{2/3} = \qty{0.16}{m^2}$.
\textbf{22.} $2.5\times 0.16\times 34 = \qty{13.6}{W}$，而假定的
$2B = \qty{11.4}{W}$：兔子须多产生一些热量，
或通过姿势和遮蔽减少散热。
\textbf{23.} 张开时：$10\times 0.01\times 34 = \qty{3.4}{W}$，又约为毛皮散热的四分之一；
闭合时：$10\times 0.01\times 3 =
\qty{0.3}{W}$。它们会闭合。
\textbf{24.} 毛皮：$2.5\times 0.16\times 9 = \qty{3.6}{W}$；耳朵张开：
$10\times 0.01\times 9 = \qty{0.9}{W}$；合计 \qty{4.5}{W}，而产热为
\qty{5.7}{W}。其余 \qty{1.2}{W} 必须通过蒸发散失：兔子喘息
（\qty{1.2}{W} 相当于每小时 \qty{1.8}{g} 水），
趴卧在凉爽的地面上，并待在阴凉处。
\textbf{25.} 每日能量消耗 \qty{0.86}{MJ}，为其\omterm{def:b1:organism-environment:allometry}{基础代谢率}的 $1.75$ 倍；
水的周转量 \qty{250}{g}，占其体重的 \qty{12.5}{\%}。
\end{solution}
